\subsection{8. Git Workflow: How and When to
Commit}\label{git-workflow-how-and-when-to-commit}

\subsubsection{The Commit Message
Convention}\label{the-commit-message-convention}

\textbf{Format:} \texttt{[EXP-NNN] action: description}

{\def\LTcaptype{none} % do not increment counter
\small
\begin{longtable}[]{@{}
  >{\raggedright\arraybackslash}p{(\linewidth - 4\tabcolsep) * \real{0.3478}}
  >{\raggedright\arraybackslash}p{(\linewidth - 4\tabcolsep) * \real{0.2609}}
  >{\raggedright\arraybackslash}p{(\linewidth - 4\tabcolsep) * \real{0.3913}}@{}}
\toprule\noalign{}
\begin{minipage}[b]{\linewidth}\raggedright
Action
\end{minipage} & \begin{minipage}[b]{\linewidth}\raggedright
When
\end{minipage} & \begin{minipage}[b]{\linewidth}\raggedright
Example
\end{minipage} \\
\midrule\noalign{}
\endhead
\bottomrule\noalign{}
\endlastfoot
\texttt{setup} & First scaffolding, cloning code &
\texttt{[EXP-001] setup: clone UBA code and create venv} \\
\texttt{run} & Running experiments, generating results
&
\texttt{[EXP-001] run: reproduce Table 3 with original params} \\
\texttt{fix} & Fixing bugs or dependency issues &
\texttt{[EXP-003] fix: adjust data path for local AEMO dataset} \\
\texttt{results} & Adding result tables or figures &
\texttt{[EXP-001] results: add comparison table for RMSE metrics} \\
\texttt{docs} & Updating README or REPLICATION\_LOG &
\texttt{[EXP-001] docs: complete REPLICATION\textbackslash\{\}\_LOG for session} \\
\texttt{refactor} & Restructuring code &
\texttt{[EXP-001] refactor: move helper functions to adapted/} \\
\end{longtable}
}

\textbf{For repo-wide changes:}

\begin{minted}{text}
[REPO] docs: update README progress summary
[REPO] feat: add shared SERA metric implementation
\end{minted}

\subsubsection{When to Commit}\label{when-to-commit}

Commit after each of these milestones: 1. After scaffolding environment
(status -\textgreater{} {[}RUNNING{]}) 2. After first successful run
(status -\textgreater{} {[}REVIEW{]}) 3. After generating each set of
results 4. After completing the experiment (status -\textgreater{}
{[}DONE{]}) 5. After any significant REPLICATION\_LOG update

\subsubsection{The Standard Commit
Sequence}\label{the-standard-commit-sequence}

\begin{minted}{bash}
# 1. Check what changed
git status

# 2. Stage only what's relevant
git add experiments/NNN_folder_name/
git add README.md  # If you ran update_readme.py

# 3. Verify what's staged
git diff --staged --stat

# 4. Commit
git commit -m "[EXP-NNN] action: short description"

# 5. Push to GitHub
git push origin main
\end{minted}

\begin{quote}
\small {[}WARN{]} \textbf{Never commit:} \texttt{*.csv} data
files, \texttt{*.h5} model files,
\texttt{.venv/} directories, or
\texttt{\textbackslash\{\}\_\textbackslash\{\}\_pycache\textbackslash\{\}\_\textbackslash\{\}\_/}. These are all in
\texttt{.gitignore}.
\end{quote}
